% =============================================================================
%  «ТЕОРИЯ СИСТЕМ» — ствол серии «Архитектура Рассуждения»
%  Блок 4. ВРЕМЯ: система, как она живёт.  Разделы 4.1–4.5 (черновик прозы).
%  Стиль: «выводим и показываем»; репо-якоря — в приложении, не в сносках.
%  Самокомпилируемо: pdflatex blok-4-vremya.tex (см. память latex-build-windows).
% =============================================================================
\documentclass[12pt,a4paper]{article}
\usepackage{cmap}
\usepackage[T2A]{fontenc}
\usepackage[utf8]{inputenc}
\usepackage[russian]{babel}
\usepackage{paratype}
\usepackage{amsmath,amssymb,amsthm}
\usepackage{listings}
\usepackage[expansion=false]{microtype}
\usepackage[margin=28mm]{geometry}

\newcommand{\rezultat}{\medskip\noindent\emph{Что отсюда следует.}\ }

\title{\textbf{Теория Систем}\\[2mm]\large Блок 4. Время: система, как она живёт}
\author{}
\date{}

\begin{document}
\maketitle
\thispagestyle{empty}

\noindent\small\textit{Черновик: разделы 4.1--4.5. Продолжает Блок~3 (категория). Статичный объект
получает время: оператор, эволюция, две стрелы. Машинные якоря --- в Приложении.}\normalsize

\bigskip

\noindent Категория показала, как системы складываются. Теперь покажем, как одна система
\emph{живёт}. Время системы --- это \emph{итерация} под её правилами; и главный результат блока
прост и неожидан: во всей динамике необратима \emph{ровно одна} стрела --- и это L5.

% =============================================================================
\section*{4.1\quad Оператор и эволюция}
% =============================================================================

Оператор \emph{внутри} системы --- это её \emph{эндо}-морфизм: система преобразует собственное
состояние, сохраняя роли. Такие операторы образуют \emph{моноид} --- есть тождество, и композиция
ассоциативна.

Оператор \emph{снаружи} --- действие из другой системы через сцепку; внутренний оператор оказывается
его частным случаем (само-сцепка). Это и есть «оператор, действующий внутри другой системы, которая
производит взаимодействие».

Эволюция --- повторное применение оператора: \emph{траектория}, \emph{равновесие} (неподвижная
точка), \emph{обратимость}. Простые примеры --- переключатель, коллапс, сдвиг.

\rezultat
Время системы --- \textbf{итерация её оператора} под правилами. Оператор внутри и снаружи --- один
механизм с разным \emph{источником} действия.

% =============================================================================
\section*{4.2\quad Две стрелы}
% =============================================================================

В динамике различимы две разные «стрелы». \emph{Состояние} возвратимо: орбита может замкнуться,
система --- вернуться в прежнее состояние. Но \emph{стрела времени} --- порядок L5 --- необратима:
пройденный шаг порядка не отменить.

Это и есть тезис блока. \textbf{Необратима ровно одна стрела --- и это L5}, а не движение состояния.
А \emph{аттрактор} --- то, к чему стягивается динамика, --- оказывается role-limit: например,
$(1/2)^n \to 0$, но ни одно $(1/2)^n$ не равно нулю.

\rezultat
Возвратность состояния и необратимость времени --- \emph{не одно и то же}. Динамика не делает время
обратимым: обратимо лишь движение по состояниям, а \textbf{стрела порядка необратима}.

% =============================================================================
\section*{4.3\quad Группа, действие, бассейн}
% =============================================================================

\emph{Обратимые} операторы: их степени тоже обратимы --- они образуют \emph{группу}. Эволюция есть
действие $\mathbb{N}$ (шаги складываются), а \emph{времена возврата} образуют под-моноид
$\mathbb{N}$: у переключателя --- чётные, у коллапса --- все, у сдвига --- только нуль.

\emph{Бассейн достижимости} --- куда вообще можно прийти. Половинное стягивание достигает нуля
\emph{только из нуля}: из любой другой точки оно приближается, но не приходит.

\rezultat
Обратимая динамика --- \textbf{группа}; времена возврата структурированы под-моноидом; а недостижимый
предел стягивания --- снова \emph{role-limit} (приближается, не приходит).

% =============================================================================
\section*{4.4\quad Сопряжённость, инварианты, Ляпунов}
% =============================================================================

\emph{Сопряжённость} --- переименование одной динамики в другую через биекцию-интертвайн ---
есть отношение эквивалентности; она сохраняет равновесия и периоды. Переключатель и коллапс
\emph{не} сопряжены: у них разные неподвижные точки.

\emph{Инвариантные} множества (замкнутые под динамикой) образуют \emph{решётку}; достижимое из старта
--- наименьшее инвариантное, его содержащее. \emph{Функция Ляпунова} (невозрастающая вдоль динамики)
держит систему ограниченной; строгая --- гарантирует, что ушедшее из равновесия не вернётся.

\rezultat
Динамика \textbf{классифицируется} (сопряжённость), \textbf{структурируется} (решётка инвариантов) и
\textbf{контролируется} (Ляпунов). Это обычная теория динамических систем --- но построенная
\emph{над триадой}, без аппарата сверх неё.

% =============================================================================
\section*{4.5\quad Аттрактор как role-limit}
% =============================================================================

Соберём звон, прошедший через весь блок. \emph{Аттрактор --- это role-limit}: динамика
\emph{приближается} к пределу, но финитно его не \emph{достигает} ($(1/2)^n \to 0$, и всё же каждое
$(1/2)^n > 0$). Это та же граница, что двойственность вырезать\,/\,слить у (ко)эквалайзеров (§3.6) и
незавершённая башня уровней (Блок~6).

\rezultat
В динамике граница теории впервые \emph{видна как предел движения}. Аттрактор существует как
\textbf{процесс-горизонт}, а не как достигаемая точка --- и это прямой мост к Блоку~7, где единственная
несущая стена будет названа.

\medskip
\noindent\emph{Переход.} Одна система во времени разобрана. Следующий блок берёт \emph{много} систем
сразу --- и спрашивает, что у целого есть сверх частей.

% =============================================================================
\appendix
\section*{Приложение к Блоку 4 — машинные якоря}
% =============================================================================
\small\raggedright\sloppy
\noindent Всё перечисленное машинно проверено; единственная аксиома ветви --- \texttt{classic} (L3).

\begin{itemize}\itemsep2pt
\item \textbf{§4.1.} оператор внутри/вне --- \texttt{InsideOperator}, \texttt{OutsideOperator},
\texttt{inside\_is\_self\_outside} (\texttt{ERROperator.v}); эволюция \texttt{evolve},
\texttt{trajectory}, \texttt{equilibrium}, \texttt{reversible}; примеры \texttt{flip}/\texttt{collapse}/\texttt{SB}
(\texttt{ERRDynamics.v}).

\item \textbf{§4.2.} две стрелы --- \texttt{arrow\_never\_returns}; аттрактор=role-limit
(\texttt{halve}/\texttt{double}).\\
\emph{Файл:} \texttt{ERRDynamicsArrow.v}.

\item \textbf{§4.3.} группа обратимых, бассейн \texttt{reaches} (\texttt{ERRDynamicsGroupBasin.v});
$\mathbb{N}$-действие \texttt{evolve\_add}, \texttt{return\_times} (\texttt{ERRDynamicsAction.v}).

\item \textbf{§4.4.} сопряжённость, \texttt{flip}$\,\not\sim\,$\texttt{collapse}
(\texttt{ERRDynamicsConjugacy.v}); решётка инвариантов, \texttt{reachable}
(\texttt{ERRDynamicsInvariant.v}); функции Ляпунова (\texttt{ERRDynamicsLyapunov.v}).

\item \textbf{§4.5.} аттрактор=role-limit --- \texttt{ERRDynamicsArrow.v}; мост к процессам ---
\texttt{ERRProcessBridge.v} (Блок~7).
\end{itemize}
\normalsize

% --- Дальше: Блок 5 (ЦЕЛОЕ СВЕРХ ЧАСТЕЙ) по тезисному плану в Книги/Теория Систем/. ---

\end{document}
